\documentclass[11pt,a4paper]{article}
\usepackage{curso}
\begin{document}
\portada{Lección}{9.3}{Modelos termodinámicos de eliminación de ruido (DTM)}

\noindent\textbf{Unidad 9.}\quad\textbf{Objetivos:} el compromiso mezcla--expresividad; DTM como cadena de EBM condicionales; el modelo de energía de la DTCA y la lectura
de la cifra de $10^4\times$.
\textbf{Prerrequisitos:} 4.3, 5.3, 9.1--9.2.

\section{Compromiso mezcla--expresividad}
Un EBM monolítico $P(x)\propto e^{-E(x)}$ que representa datos multimodales tiene valles profundos separados por barreras; un muestreador local las cruza con probabilidad
$\sim e^{-\Delta E}$: más expresividad, mezcla más lenta.
\begin{ejemplo}
\evS\ Objetivo $0{,}8\,\N(2;0{,}3^2)+0{,}2\,\N(-2;0{,}3^2)$; Langevin desde $\N(0,1)$, 8000 pasos: fracción en el modo derecho $51\,\%$ (debería ser $80\,\%$).
\end{ejemplo}

\section{DTM}
Como en la difusión, se añade ruido en $T$ pasos y cada paso inverso es un EBM condicional:
\[P_\theta(x_{t-1}\mid x_t)\propto e^{-E^\theta_{t-1}(x_{t-1})}\,q(x_t\mid x_{t-1}),\]
casi unimodal y fácil de muestrear; la complejidad se construye a lo largo de la cadena.
\begin{ejemplo}
\evS\ Mismo objetivo, 32 niveles con 40 pasos de Langevin cada uno: $75\,\%$ en el modo derecho. El error restante proviene de condicionales aún bimodales y de la discretización;
con muy pocos niveles, los condicionales vuelven a ser multimodales.
\end{ejemplo}

\section{Arquitectura DTCA}
Cada condicional es una máquina de Boltzmann dispersa: rejilla $L\times L$ ($L=70$) de variables binarias, grado 12, nodos de datos y latentes. Cada nodo es un generador de bits
aleatorios solo de transistores con respuesta sigmoidal medida y $\tau_0\approx100$ ns \evM. Rejilla bipartita: un barrido $\approx2\tau_0$.

\section{Energía}
\[E=T\,K_{\text{mix}}\,L^2E_{\text{cell}},\qquad E_{\text{cell}}=E_{\text{rng}}+E_{\text{bias}}+E_{\text{clock}}+E_{\text{comm}}\approx2\ \text{fJ}.\]
Con $T=8$, $K_{\text{mix}}=250$, $L=70$: $\approx20$ nJ por muestra y $T K_{\text{mix}}2\tau_0\approx0{,}4$ ms. Los autores estiman igualar el FID de modelos en GPU en Fashion-MNIST
binarizado con $\sim10^4$ veces menos energía \evE. Advertencias: combina mediciones del RNG con modelos y simulaciones; la energía de GPU se calcula como FLOPs/especificación
(subestimada); banco pequeño. Híbrido DTM + red neuronal: paridad en CIFAR-10 con una red diez veces menor.

\begin{aplicacion}
Modelado generativo en hardware probabilístico; penalización adaptativa de correlación (ACP) para mantener la mezcla durante el entrenamiento. Laboratorio L14: DTM en THRML
y su energía con la fórmula de la DTCA.
\end{aplicacion}

\begin{error}
\begin{itemize}[nosep]
\item EBM monolítico con cadenas cortas. \quad Muy pocos pasos de eliminación de ruido.
\item Confundir $\tau_0$ con $K_{\text{mix}}$. \quad Citar $10^4\times$ como medición de un chip.
\end{itemize}
\end{error}

\begin{thebibliography}{9}\small
\bibitem[48]{dtm} A. Jelinčič et al., An efficient probabilistic hardware architecture for diffusion-like models, arXiv:2510.23972 (2025).
\bibitem[22]{z1} Extropic, From One to One Billion: Torx, Thermalizers, and Z1, 2026.
\bibitem{ho} J. Ho, A. Jain y P. Abbeel, Denoising diffusion probabilistic models, \emph{NeurIPS} (2020).
\end{thebibliography}
\end{document}
